% !TEX program = xelatex
% Document généré par Lean AI Station. Compiler avec XeLaTeX (sur Overleaf : Menu → Compilateur → XeLaTeX).
\documentclass[11pt]{article}
\usepackage[a4paper,margin=2.5cm]{geometry}
\usepackage{fontspec}
\setmonofont{DejaVu Sans Mono}[Scale=MatchLowercase]
\usepackage[french]{babel}
\usepackage{amsmath,amssymb,amsthm}
\usepackage{fvextra}
\usepackage[dvipsnames]{xcolor}
\theoremstyle{plain}
\newtheorem{theorem}{Théorème}

\title{Énoncé et preuve vérifiée par Lean}
\author{}
\date{}

\begin{document}
\maketitle

\section*{Énoncé}
\begin{theorem}
Montrer que la somme de deux entiers pairs est paire.
\end{theorem}

\section*{Énoncé formalisé en Lean 4}
\begin{Verbatim}[breaklines=true,frame=single,fontsize=\small,rulecolor=\color{black!25}]
theorem mon_probleme : ∀ a b : ℤ, Even a → Even b → Even (a + b) := by sorry
\end{Verbatim}

\section*{Preuve formelle}
Preuve vérifiée par Lean (Prouveur, Lean 4.9.0-rc1) le 04/10/2026.
\begin{Verbatim}[breaklines=true,frame=single,fontsize=\small,rulecolor=\color{black!25}]
import Mathlib
import Aesop

set_option maxHeartbeats 400000

open BigOperators Real Nat Topology Rat

theorem mon_probleme : ∀ a b : ℤ, Even a → Even b → Even (a + b) := by
  have h_main : ∀ a b : ℤ, Even a → Even b → Even (a + b) := by
    intro a b h_even_a h_even_b
    
    cases' h_even_a with k hk
    cases' h_even_b with m hm
    
    use k + m
    
    <;> simp [hk, hm, Int.mul_add, Int.add_mul]
    <;> ring
    <;> omega
  exact h_main
\end{Verbatim}

\end{document}
